\documentclass[10pt,aspectratio=169]{beamer}

\usepackage[utf8]{inputenc}
\usepackage[spanish]{babel}
\usepackage[T1]{fontenc}
\usepackage{booktabs}
\usepackage{array}
\usepackage{tikz}
\usetikzlibrary{positioning,arrows.meta}

\definecolor{azulprofundo}{HTML}{1e3a5f}
\definecolor{verdeecobrief}{HTML}{16a34a}
\definecolor{gristexto}{HTML}{334155}
\definecolor{grisclaro}{HTML}{f8fafc}

\usetheme{Madrid}
\usecolortheme{default}

\setbeamercolor{normal text}{fg=gristexto,bg=white}
\setbeamercolor{structure}{fg=verdeecobrief}
\setbeamercolor{frametitle}{fg=white,bg=azulprofundo}
\setbeamercolor{title}{fg=azulprofundo}
\setbeamercolor{subtitle}{fg=gristexto}
\setbeamercolor{block title}{fg=white,bg=azulprofundo}
\setbeamercolor{block body}{fg=gristexto,bg=grisclaro}
\setbeamercolor{alerted text}{fg=verdeecobrief}
\setbeamercolor{palette primary}{fg=white,bg=azulprofundo}
\setbeamercolor{palette secondary}{fg=white,bg=azulprofundo}
\setbeamercolor{palette tertiary}{fg=white,bg=azulprofundo}
\setbeamercolor{palette quaternary}{fg=white,bg=azulprofundo}

\setbeamertemplate{navigation symbols}{}
\setbeamertemplate{itemize item}{\color{verdeecobrief}\small$\blacktriangleright$}
\setbeamertemplate{itemize subitem}{\color{verdeecobrief}\scriptsize$\blacktriangleright$}
\setbeamertemplate{enumerate item}{\color{verdeecobrief}\insertenumlabel.}
\setbeamertemplate{footline}[frame number]

\newcommand{\headline}[1]{%
    \begin{center}
    {\Large\bfseries\color{azulprofundo} #1}
    \end{center}
}

\newcommand{\keymessage}[1]{%
    \begin{center}
    \fbox{\begin{minipage}{0.86\textwidth}
    \centering\Large\bfseries\color{verdeecobrief} #1
    \end{minipage}}
    \end{center}
}

\title{EcoBrief Bolivia}
\subtitle{IA responsable para reducir el desperdicio digital informativo}
\author{Josoe Ichuta}
\institute{AI Battle}
\date{2026}

\begin{document}

\begin{frame}
    \titlepage
    \vspace{-0.2cm}
    \keymessage{Menos ruido. Mas informacion esencial.}
\end{frame}

\begin{frame}{SITUATION}
    \headline{Queremos estar informados}
    \vspace{0.2cm}
    \begin{columns}[T]
        \column{0.48\textwidth}
        \begin{block}{Fuentes habituales}
            \begin{itemize}
                \item TV y noticieros
                \item Radio
                \item Sitios web
                \item Redes sociales
                \item Multiples titulares
            \end{itemize}
        \end{block}
        \column{0.48\textwidth}
        \begin{block}{Efecto}
            \Large
            Informarse toma tiempo porque el flujo esta fragmentado.
        \end{block}
    \end{columns}
\end{frame}

\begin{frame}{SITUATION}
    \headline{El costo oculto de informarse}
    \vspace{0.2cm}
    \begin{table}
        \centering
        \begin{tabular}{ll}
            \toprule
            \textbf{Desperdicio} & \textbf{Ejemplo} \\
            \midrule
            Tiempo & Comparar titulares repetidos \\
            Datos & Cargar paginas completas \\
            Atencion & Scrollear sin fin \\
            Confianza & Contenido sin fuente clara \\
            \bottomrule
        \end{tabular}
    \end{table}
    \vspace{0.2cm}
    \keymessage{No solo consumimos noticias: consumimos tiempo, datos y atencion.}
\end{frame}

\begin{frame}{SITUATION}
    \headline{La pregunta}
    \vspace{0.5cm}
    \keymessage{Como reducimos el ruido sin dejar de informarnos?}
    \vspace{0.5cm}
    \begin{center}
        \large Una solucion util no debe crear mas contenido.\\
        Debe reducir el contenido innecesario.
    \end{center}
\end{frame}

\begin{frame}{TASK}
    \headline{Cuatro piezas necesarias}
    \begin{enumerate}
        \item Recolectar noticias desde fuentes bolivianas.
        \item Filtrar, deduplicar, clasificar y rankear.
        \item Presentar briefs relevantes al usuario.
        \item Ejecutar el proceso varias veces al dia.
    \end{enumerate}
    \vspace{0.4cm}
    \keymessage{La IA es el ultimo paso, no el primero.}
\end{frame}

\begin{frame}{ACTION}
    \headline{Arquitectura del sistema}
    \centering
    \begin{tikzpicture}[
        node distance=0.9cm and 1.2cm,
        box/.style={rectangle, rounded corners, draw=azulprofundo, thick, fill=grisclaro, minimum width=2.8cm, minimum height=0.8cm, align=center},
        arrow/.style={-{Latex}, thick, color=verdeecobrief}
    ]
        \node[box] (scrapers) {Scrapers};
        \node[box, right=of scrapers] (api) {FastAPI};
        \node[box, right=of api] (front) {React + Vite};
        \node[box, below=of api] (db) {PostgreSQL};
        \node[box, above=of api] (llm) {LLM Provider};
        \node[box, above=of front] (providers) {Groq / OpenAI / GitHub};
        \node[box, below=of front] (dist) {Email / Telegram / WhatsApp};

        \draw[arrow] (scrapers) -- (api);
        \draw[arrow] (api) -- (front);
        \draw[arrow] (api) -- (db);
        \draw[arrow] (api) -- (llm);
        \draw[arrow] (llm) -- (providers);
        \draw[arrow] (api) -- (dist);
    \end{tikzpicture}
    \vspace{0.2cm}
    \begin{center}
        Python 3.11 \quad | \quad FastAPI \quad | \quad PostgreSQL \quad | \quad React \quad | \quad Docker
    \end{center}
\end{frame}

\begin{frame}{ACTION}
    \headline{Pipeline antes de IA}
    \centering
    \begin{tikzpicture}[
        node distance=0.45cm,
        pbox/.style={rectangle, rounded corners, draw=azulprofundo, thick, fill=grisclaro, minimum width=6.8cm, minimum height=0.52cm, align=center},
        arrow/.style={-{Latex}, thick, color=verdeecobrief}
    ]
        \node[pbox] (a) {Fuentes de noticias};
        \node[pbox, below=of a] (b) {Scraping};
        \node[pbox, below=of b] (c) {Filtro de calidad};
        \node[pbox, below=of c] (d) {Deduplicacion};
        \node[pbox, below=of d] (e) {Clasificacion};
        \node[pbox, below=of e] (f) {Ranking};
        \node[pbox, below=of f] (g) {Seleccion};
        \node[pbox, below=of g] (h) {IA para resumen};
        \draw[arrow] (a) -- (b);
        \draw[arrow] (b) -- (c);
        \draw[arrow] (c) -- (d);
        \draw[arrow] (d) -- (e);
        \draw[arrow] (e) -- (f);
        \draw[arrow] (f) -- (g);
        \draw[arrow] (g) -- (h);
    \end{tikzpicture}
    \vspace{0.15cm}
    \keymessage{No resumimos todo. Resumimos lo mas relevante.}
\end{frame}

\begin{frame}{ACTION}
    \headline{Menos ruido antes del modelo}
    \begin{columns}[T]
        \column{0.48\textwidth}
        \begin{block}{Deduplicacion}
            \begin{itemize}
                \item URL hash
                \item Titulo normalizado
                \item Huella historica SHA-256
            \end{itemize}
        \end{block}
        \column{0.48\textwidth}
        \begin{block}{Clasificacion}
            \begin{itemize}
                \item Economia
                \item Politica
                \item Deportes
                \item Tecnologia
                \item Entretenimiento
            \end{itemize}
        \end{block}
    \end{columns}
    \vspace{0.25cm}
    \begin{center}
        \large Reglas ponderadas por titulo, descripcion, contenido y categoria de fuente.
    \end{center}
\end{frame}

\begin{frame}{ACTION}
    \headline{Score de prioridad}
    \begin{table}
        \centering
        \begin{tabular}{lr}
            \toprule
            \textbf{Factor} & \textbf{Peso} \\
            \midrule
            Relevancia local & 20\% \\
            Impacto informativo & 20\% \\
            Calidad del contenido & 17\% \\
            Actualidad & 15\% \\
            Fuente & 10\% \\
            Corroboracion & 10\% \\
            Confianza de categoria & 8\% \\
            \bottomrule
        \end{tabular}
    \end{table}
    \keymessage{Cada noticia recibe un puntaje de 0 a 100.}
\end{frame}

\begin{frame}{ACTION}
    \headline{IA solo cuando aporta valor}
    \vspace{0.2cm}
    \keymessage{Filtrado -> Deduplicado -> Rankeado -> Resumido}
    \vspace{0.25cm}
    \begin{columns}[T]
        \column{0.32\textwidth}
        \begin{block}{Groq}
            Rapido y costo bajo
        \end{block}
        \column{0.32\textwidth}
        \begin{block}{OpenAI}
            Mayor capacidad
        \end{block}
        \column{0.32\textwidth}
        \begin{block}{GitHub Models}
            Alternativa con cuota
        \end{block}
    \end{columns}
    \vspace{0.3cm}
    \begin{center}
        Cada articulo filtrado antes de IA es una llamada menos al modelo.
    \end{center}
\end{frame}

\begin{frame}{RESULT}
    \headline{Producto desplegado}
    \begin{columns}[T]
        \column{0.52\textwidth}
        \begin{itemize}
            \item Scraping real de fuentes bolivianas.
            \item Resumenes generados con IA.
            \item Frontend web con noticias destacadas.
            \item Noticias recolectadas con trazabilidad.
        \end{itemize}
        \column{0.44\textwidth}
        \begin{itemize}
            \item Pagina de datos.
            \item Panel de impacto.
            \item Suscripcion por preferencias.
            \item Pruebas automatizadas.
        \end{itemize}
    \end{columns}
    \vspace{0.35cm}
    \keymessage{https://briefs.ecobriefbolivia.online/}
\end{frame}

\begin{frame}{RESULT}
    \headline{Resultado de una corrida real}
    \begin{table}
        \centering
        \begin{tabular}{lr}
            \toprule
            \textbf{Metrica} & \textbf{Valor} \\
            \midrule
            Noticias recolectadas & 174 \\
            Articulos utiles & 173 \\
            Articulos unicos & 169 \\
            Briefs generados & 39 \\
            Paginas evitadas & 135 \\
            Minutos ahorrados & 67.5 \\
            Reduccion & 77.6\% \\
            \bottomrule
        \end{tabular}
    \end{table}
    \begin{center}
        \small Corrida real: 15 de junio de 2026
    \end{center}
\end{frame}

\begin{frame}{RESULT}
    \headline{Recorrido del producto}
    \begin{enumerate}
        \item Home: briefs principales.
        \item Noticias recolectadas: fuente, imagen y fecha.
        \item Datos: indicadores externos.
        \item Impacto: metricas globales y por corrida.
        \item Suscripcion: categorias, frecuencia y canal.
    \end{enumerate}
    \vspace{0.35cm}
    \keymessage{Briefs trazables, metricas visibles y preferencias por usuario.}
\end{frame}

\begin{frame}{RESULT}
    \headline{Distribucion personalizada}
    \begin{table}
        \centering
        \begin{tabular}{ll}
            \toprule
            \textbf{Canal} & \textbf{Estado} \\
            \midrule
            Email & Funcional, pausado por cuenta personal \\
            Telegram & En integracion \\
            WhatsApp & Planteado con Twilio, pausado por costo \\
            \bottomrule
        \end{tabular}
    \end{table}
    \vspace{0.25cm}
    \begin{center}
        Corto plazo: metricas y analiticas \quad | \quad Mediano plazo: canales en produccion
    \end{center}
\end{frame}

\begin{frame}{CIERRE}
    \headline{Menos ruido. Mas informacion esencial.}
    \begin{itemize}
        \item Menos paginas abiertas.
        \item Menos duplicidad.
        \item Menos llamadas innecesarias a IA.
        \item Mas informacion trazable.
    \end{itemize}
    \vspace{0.5cm}
    \keymessage{EcoBrief Bolivia usa IA no para producir mas ruido, sino para reducirlo.}
\end{frame}

\end{document}

